% =========================================================
%  TEAQUEST — PROJECT REPORT
%  Compile with pdflatex (or upload to Overleaf)
% =========================================================

\documentclass[11pt,a4paper]{article}

% ---------------------------------------------------------
%  PACKAGES
% ---------------------------------------------------------

\usepackage[a4paper,margin=2.5cm]{geometry}
\usepackage[T1]{fontenc}
\usepackage[utf8]{inputenc}
\usepackage{lmodern}
\usepackage{xcolor}
\usepackage{graphicx}
\usepackage{booktabs}
\usepackage{longtable}
\usepackage{array}
\usepackage{enumitem}
\usepackage{fancyhdr}
\usepackage{titlesec}
\usepackage{caption}
\usepackage{listings}
\usepackage{hyperref}

% ---------------------------------------------------------
%  THEME COLORS (matching the TeaQuest palette)
% ---------------------------------------------------------

\definecolor{teagreen}{RGB}{117,214,107}
\definecolor{teadark}{RGB}{17,29,23}
\definecolor{teacream}{RGB}{244,230,189}
\definecolor{teagold}{RGB}{245,201,93}
\definecolor{codegreen}{RGB}{80,160,80}
\definecolor{codegray}{RGB}{120,120,120}

\hypersetup{
    colorlinks=true,
    linkcolor=teadark,
    urlcolor=teadark,
    citecolor=teadark,
    pdftitle={TeaQuest Project Report},
    pdfauthor={TeaQuest Team}
}

% ---------------------------------------------------------
%  LISTINGS STYLE
% ---------------------------------------------------------

\lstdefinestyle{codeStyle}{
    backgroundcolor=\color{teadark!4},
    basicstyle=\ttfamily\footnotesize,
    keywordstyle=\color{codegreen}\bfseries,
    commentstyle=\color{codegray}\itshape,
    stringstyle=\color{teadark},
    numbers=left,
    numberstyle=\tiny\color{codegray},
    numbersep=8pt,
    breaklines=true,
    frame=single,
    rulecolor=\color{teadark!40},
    framesep=6pt,
    tabsize=4,
    showstringspaces=false,
    captionpos=b
}

\lstset{style=codeStyle}

% ---------------------------------------------------------
%  HEADINGS & HEADER/FOOTER
% ---------------------------------------------------------

\titleformat{\section}
  {\Large\bfseries\color{teadark}}{\thesection}{0.8em}{}
\titleformat{\subsection}
  {\large\bfseries\color{teadark}}{\thesubsection}{0.8em}{}

\pagestyle{fancy}
\fancyhf{}
\fancyhead[L]{\small\textsc{TeaQuest --- Project Report}}
\fancyhead[R]{\small\thepage}
\renewcommand{\headrulewidth}{0.4pt}

% ---------------------------------------------------------
%  TITLE PAGE DATA
% ---------------------------------------------------------

\newcommand{\ReportTitle}{TEAQUEST}
\newcommand{\ReportSubtitle}{Brew Your Adventure --- A Pixel-Art Fantasy E-Commerce Platform}
\newcommand{\ReportAuthor}{[Your Name]}
\newcommand{\ReportDate}{August 2026}

\begin{document}

% =========================================================
%  TITLE PAGE
% =========================================================

\begin{titlepage}
    \centering
    \vspace*{2cm}

    {\color{teagreen}\rule{\linewidth}{2pt}}\\[0.6cm]

    {\Huge\bfseries\scshape \ReportTitle\par}
    \vspace{0.5cm}
    {\color{teagreen}\rule{12pt}{12pt}}%
    {\color{teagold}\rule{12pt}{12pt}}%
    {\color{teagreen}\rule{12pt}{12pt}}\par
    \vspace{0.5cm}
    {\Large\itshape \ReportSubtitle\par}

    \vspace{0.5cm}
    {\color{teagreen}\rule{\linewidth}{2pt}}\\[2cm]

    {\large Project Technical Report\par}
    \vspace{1.5cm}

    \begin{tabular}{rl}
        \textbf{Author}      & \ReportAuthor \\[4pt]
        \textbf{Version}     & 1.0 \\[4pt]
        \textbf{Date}        & \ReportDate \\[4pt]
        \textbf{Live Demo}   & \url{https://riyadh5674.github.io/TeaQuest/} \\
    \end{tabular}

    \vfill
    {\small\itshape ``Discover. Brew. Explore.''\par}
    \vspace{1cm}
\end{titlepage}

% =========================================================
%  ABSTRACT
% =========================================================

\begin{abstract}
\noindent
\textbf{TeaQuest} is a premium tea e-commerce web application wrapped in an
original pixel-art fantasy theme. The project combines a complete shopping
experience --- browsing, filtering, favorites, cart, checkout and order
history --- with game-inspired mechanics such as experience points, quests,
achievements, a discovery codex, a randomized ``Tea Roulette'' and a
personality-based ``Tea Oracle'' recommender. The frontend is built entirely
with vanilla HTML5, CSS3 and JavaScript, requiring no frameworks or build
tools. Persistence uses a hybrid model: a local-first \texttt{localStorage}
layer with a Supabase backend (PostgreSQL, authentication and row-level
security) layered on top. This report documents the system architecture,
frontend design, functional modules, database schema, security model and
known limitations, and outlines the roadmap toward a production-ready
platform.
\end{abstract}

\vspace{0.5cm}
{\color{teadark!30}\hrule}
\vspace{0.5cm}

\tableofcontents
\newpage

% =========================================================
%  1. INTRODUCTION
% =========================================================

\section{Introduction}

Online tea shopping is functionally identical across most stores: a grid of
products, a cart and a checkout form. \textbf{TeaQuest} challenges this
convention by framing the entire purchase journey as a cozy role-playing
game. Customers are ``players'', products are ``loot'', orders are
``quests'', and the administrator is the ``Guild Master''. The goal is
simple: \emph{make buying tea easy while making the experience memorable}.

The project is implemented as a single-page application (SPA) with no
frameworks, no package installation and no build step. Despite this
deliberately minimal toolchain, it delivers features typically associated
with much larger stacks: persistent user accounts, an admin dashboard with
analytics, server-side authorization rules and structured cloud data.

\subsection{Motivation}

\begin{itemize}[nosep]
    \item Demonstrate that strong gamification can coexist with serious
          e-commerce usability.
    \item Prove that a modern full-stack application can be delivered with
          vanilla web technologies plus a Backend-as-a-Service provider.
    \item Create an original visual identity that evokes retro games without
          using any copyrighted assets.
\end{itemize}

\subsection{Scope}

The current version covers the complete storefront, customer accounts, an
administrative back office and a cloud database with enforced security
policies. Payment processing is intentionally out of scope and simulated;
real payments are planned for a future release.

% =========================================================
%  2. OBJECTIVES
% =========================================================

\section{Objectives}

\begin{enumerate}[nosep]
    \item Build a responsive storefront with search, category filtering,
          sorting and product detail views.
    \item Implement cart, favorites and a multi-step checkout flow with
          order history.
    \item Add a gamification layer: XP and levels, quests, achievements,
          a discovery codex, Tea Roulette and the Tea Oracle quiz.
    \item Provide customer registration/login with persistent profiles.
    \item Provide an admin dashboard for products, orders, customers and
          settings including JSON backup/restore.
    \item Persist all state in \texttt{localStorage} with an optional
          Supabase cloud backend protected by row-level security (RLS).
    \item Keep the project dependency-free: no frameworks, no bundlers,
          no npm install.
\end{enumerate}

% =========================================================
%  3. TECHNOLOGY STACK
% =========================================================

\section{Technology Stack}

\begin{table}[h!]
\centering
\caption{Technologies and libraries used.}
\begin{tabular}{@{}p{4.2cm}p{4.2cm}p{6.2cm}@{}}
\toprule
\textbf{Layer} & \textbf{Technology} & \textbf{Purpose} \\
\midrule
Markup       & HTML5                  & Single-file SPA shell, semantic sections and modals \\
Styling      & CSS3                   & Pixel-art theme, animations, responsive layout, theming via custom properties \\
Logic        & Vanilla JavaScript (ES2020+) & Application engine, rendering, state management \\
Persistence  & \texttt{localStorage}  & Offline-first data store for products, users, cart, favorites, orders \\
Backend      & Supabase (PostgreSQL)  & Cloud database, authentication, row-level security \\
Client SDK   & supabase-js v2 (CDN)   & Database and auth access from the browser \\
Fonts        & Press Start 2P, VT323  & Retro pixel typography via Google Fonts \\
\bottomrule
\end{tabular}
\end{table}

No Node.js, npm or bundler is required; the site runs by opening
\texttt{index.html} through any static server (e.g., VS Code Live Server).

% =========================================================
%  4. SYSTEM ARCHITECTURE
% =========================================================

\section{System Architecture}

TeaQuest follows a \textbf{local-first, cloud-enhanced} architecture. All
features work fully offline against \texttt{localStorage}; when the Supabase
SDK is reachable, accounts and profiles are transparently synchronized to
the cloud.

\subsection{High-Level Overview}

\begin{lstlisting}[caption={Layered architecture (simplified).},language={}]
+----------------------------------------------------------+
|                      BROWSER (SPA)                        |
|                                                          |
|  index.html ..... pages, modals, drawers (static shell)   |
|  java.js ........ application engine (~97 functions)      |
|  style.css ...... pixel-art presentation layer            |
|                                                          |
|  State: products | users | cart | favorites | orders      |
|         currentUser                                       |
+---------------------------+------------------------------+
                            |  sync (profiles, auth)
                            v
+----------------------------------------------------------+
|                    SUPABASE (cloud)                       |
|  Auth ......... email/password sessions                   |
|  PostgreSQL ... profiles, products, orders,               |
|                 order_items, favorites                    |
|  RLS .......... policy-enforced authorization             |
+----------------------------------------------------------+
\end{lstlisting}

\subsection{State Management}

Application state lives in module-level variables initialized from storage:

\begin{itemize}[nosep]
    \item \texttt{products}, \texttt{users}, \texttt{cart},
          \texttt{favorites}, \texttt{orders}, \texttt{currentUser}
    \item UI state: \texttt{activeCategory}, \texttt{authMode}
\end{itemize}

Every mutation follows the same pattern: \emph{modify state $\rightarrow$
persist via \texttt{saveStorage()} $\rightarrow$ re-render affected views}.
A single \texttt{renderEverything()} routine refreshes all major surfaces
after bootstrapping.

\subsection{Boot Sequence}

On \texttt{DOMContentLoaded} the engine runs each feature initializer
(navigation, theme, shop, cart, favorites, authentication, checkout,
contact, admin, roulette, oracle), then restores the Supabase session via
\texttt{restoreSession()} and performs the initial full render.

% =========================================================
%  5. FRONTEND DESIGN
% =========================================================

\section{Frontend Design and User Experience}

\subsection{Visual Identity}

The design language is an original \emph{cozy pixel fantasy}: deep forest
greens, cream and gold accents, hard offset shadows that mimic chunky pixel
borders, scanline atmosphere, drifting clouds, steam particles and a
hand-built CSS tea house on the hero section. Two Google Fonts carry the
identity: \textbf{Press Start 2P} for headings and \textbf{VT323} for body
text.

\subsection{Single-Page Navigation}

All pages are \texttt{<section>} elements toggled through a hash-aware
router (\texttt{showPage()}), keeping browser back/forward functional
without a framework. Pages: \textbf{Home, Shop, Codex, Lore, Contact,
Profile, Admin}. Overlays include the cart drawer, favorites drawer and
seven modals (auth, product details, checkout, admin product editor,
roulette, oracle, order details).

\subsection{Accessibility and Motion}

\begin{itemize}[nosep]
    \item Visible focus outlines (\texttt{:focus-visible}) on all
          interactive elements.
    \item Full \texttt{prefers-reduced-motion} support disabling
          animations and smooth scrolling.
    \item Responsive layout with a dedicated mobile navigation drawer.
\end{itemize}

\subsection{Theming}

A light/dark ``world'' toggle swaps CSS custom properties at runtime and
persists the choice, letting players switch between day and night realms.

% =========================================================
%  6. FUNCTIONAL MODULES
% =========================================================

\section{Functional Modules}

\subsection{E-Commerce Engine}

\begin{itemize}[nosep]
    \item \textbf{Catalog}: eight seeded teas across four categories
          (green, black, herbal, special) with rarity tiers from
          \emph{common} to \emph{legendary}.
    \item \textbf{Discovery}: live text search, category filters and
          sorting (featured, price ascending/descending, rating).
    \item \textbf{Cart}: quantity management, stale-item pruning against
          the current catalog, running total and slide-in drawer.
    \item \textbf{Favorites}: heart-toggle per product with a dedicated
          drawer and badge counter.
    \item \textbf{Checkout}: summary review, delivery details form,
          payment method selection (cash on delivery or demo card) and
          order placement.
    \item \textbf{Order history}: per-customer quest history with status
          tracking.
\end{itemize}

\subsection{Gamification Layer}

\begin{table}[h!]
\centering
\caption{Game systems and their triggers.}
\begin{tabular}{@{}p{3.4cm}p{5.6cm}p{5.6cm}@{}}
\toprule
\textbf{System} & \textbf{Description} & \textbf{Progress Source} \\
\midrule
XP \& Levels    & Points awarded for discoveries and actions; animated level-up overlay & \texttt{user.xp} \\
Quests (5)      & Goals such as FIRST BREW, RELIC HUNTER, LUCKY LEAF & Discoveries, spins, completed orders \\
Achievements (6)& Permanent badges, e.g.\ TEA MASTER (Level 10), CODEX KEEPER (all teas) & Profile milestones \\
Tea Codex       & Collection log recording every discovered tea forever & \texttt{user.discoveries} \\
Tea Roulette    & Slot-machine style random tea picker with spin counter & \texttt{roulette\_spins} \\
Tea Oracle      & Mood/flavor/strength quiz returning a matched brew & Product metadata matching \\
Easter Egg      & Konami code unlocks a secret hue-shifted realm & Keyboard sequence listener \\
\bottomrule
\end{tabular}
\end{table}

\subsection{Authentication and Profiles}

Accounts use Supabase email/password auth. On first login a profile row is
auto-created (database trigger) and mapped into the app shape
(\texttt{mapProfile}). XP, discoveries, achievements and roulette spins are
written back to the cloud on every change via \texttt{syncProfile()}. If
Supabase is unavailable the app degrades gracefully to local-only mode.

For the offline fallback path, passwords are hashed with SHA-256
(Web Crypto) using a static application salt, with an FNV-1a fallback for
environments without \texttt{crypto.subtle}.

\subsection{Admin Dashboard (``Guild Master Mode'')}

\begin{itemize}[nosep]
    \item \textbf{Overview}: revenue, order, player and product counters;
          a seven-day revenue chart; recent-order feed.
    \item \textbf{Products}: searchable table with create/edit/delete via
          a modal editor.
    \item \textbf{Orders}: status filters (PROCESSING, SHIPPED, DELIVERED,
          CANCELLED), inline status updates, detail view and deletion.
    \item \textbf{Customers}: player list with level, order count, role
          toggle and removal.
    \item \textbf{Settings}: password change, JSON export/import backup and
          a full factory reset.
\end{itemize}

% =========================================================
%  7. DATABASE DESIGN
% =========================================================

\section{Database Design (Supabase)}

The schema (\texttt{supabase/schema.sql}) defines five tables. Products use
readable text identifiers (\texttt{tea-001} \ldots \texttt{tea-008})
matching the frontend seed data exactly.

\begin{table}[h!]
\centering
\caption{Database tables.}
\begin{tabular}{@{}p{3.2cm}p{4.6cm}p{6.8cm}@{}}
\toprule
\textbf{Table} & \textbf{Key Columns} & \textbf{Notes} \\
\midrule
\texttt{profiles}     & id (FK auth.users), name, email, role, xp, discoveries, achievements, roulette\_spins & One row per player; \texttt{role} restricted to \emph{player}/\emph{admin} \\
\texttt{products}     & id (text PK), name, category, price, icon, rating, rarity, origin, flavor\_notes, moods, flavor\_profile, strength, stock & Seeded with the eight launch teas; idempotent seed \\
\texttt{orders}       & id (uuid PK), user\_id FK, customer, address, phone, payment, total, status & Status constrained to four workflow states \\
\texttt{order\_items} & order\_id FK, product\_id FK, name, price, quantity & Line items cascade with their order \\
\texttt{favorites}    & (user\_id, product\_id) composite PK & Private per-player saved loot \\
\bottomrule
\end{tabular}
\end{table}

\subsection{Automation Triggers}

\begin{itemize}[nosep]
    \item \texttt{on\_auth\_user\_created} --- inserts a \texttt{profiles}
          row whenever a user signs up, deriving the display name from
          metadata or the email prefix.
    \item \texttt{protect\_profile\_role} --- blocks non-admins from
          changing \texttt{profiles.role}, preventing privilege
          escalation through the public API.
\end{itemize}

\subsection{Representative Policy Excerpt}

\begin{lstlisting}[language=SQL,caption={Row-level security: products are world-readable, admin-writable.}]
create policy "products_select"
    on public.products for select
    using (true);

create policy "products_insert"
    on public.products for insert
    with check (public.is_admin ());
\end{lstlisting}

% =========================================================
%  8. SECURITY MODEL
% =========================================================

\section{Security Model}

\subsection{Server-Side Authorization}

Authorization is enforced \emph{inside the database}, not in JavaScript.
Every table enables RLS; a \texttt{security definer} helper
(\texttt{is\_admin()}) resolves the caller's role without exposing the
profiles table to enumeration.

\begin{itemize}[nosep]
    \item Players read/update only their own profile; admins see all.
    \item Orders are visible and creatable only by their owner; status
          transitions are additionally manageable by admins.
    \item Order items inherit visibility from the parent order.
    \item Favorites are strictly private per user.
    \item Role changes require admin privileges (trigger-enforced).
\end{itemize}

\subsection{Known Findings and Recommendations}

\begin{enumerate}[nosep]
    \item \textbf{Client-side admin credential.} A fallback admin password
          is embedded in \texttt{java.js}. Although real authorization is
          enforced by RLS, the constant should be removed so no usable
          credential ships in public source.
    \item \textbf{Legacy demo account.} The offline \texttt{defaultUsers}
          seed still contains a plaintext demo password; it should be
          dropped now that Supabase auth is primary.
    \item \textbf{Static hashing salt.} The local-fallback hash uses one
          fixed salt; acceptable for a demo path, but the long-term answer
          is Supabase-managed auth exclusively.
    \item \textbf{Documentation drift.} The README still describes the
          pre-backend prototype; it should be updated to reflect the
          Supabase integration.
\end{enumerate}

% =========================================================
%  9. PROJECT METRICS
% =========================================================

\section{Project Metrics}

\begin{table}[h!]
\centering
\caption{Codebase statistics (current working tree).}
\begin{tabular}{@{}lrrl@{}}
\toprule
\textbf{File} & \textbf{Lines} & \textbf{Size (KB)} & \textbf{Role} \\
\midrule
index.html            & 1{,}609 & $\approx$62  & SPA shell, 7 pages, 8 overlays \\
style.css             & 5{,}142 & $\approx$69  & Theme, animation, responsiveness \\
java.js               & 5{,}803 & $\approx$107 & Engine, $\sim$97 functions \\
supabase/schema.sql   & 371     & $\approx$10  & Tables, RLS, triggers, seed \\
\midrule
\textbf{Total}        & \textbf{12{,}925} & $\approx$\textbf{248} & \\
\bottomrule
\end{tabular}
\end{table}

\subsection{Seeded Product Catalog}

\begin{table}[h!]
\centering
\caption{Launch lineup (eight teas).}
\begin{tabular}{@{}llllr l@{}}
\toprule
\textbf{ID} & \textbf{Name} & \textbf{Category} & \textbf{Rarity} & \textbf{Price} & \textbf{Rating} \\
\midrule
tea-001 & Forest Mint    & green   & common    & \$12 & 4.9 \\
tea-002 & Dragon Ember   & black   & uncommon  & \$15 & 4.8 \\
tea-003 & Moon Blossom   & herbal  & uncommon  & \$14 & 5.0 \\
tea-004 & Emerald Mist   & green   & rare      & \$18 & 4.9 \\
tea-005 & Sunset Chai    & special & rare      & \$17 & 4.7 \\
tea-006 & Wizard's Berry & herbal  & rare      & \$16 & 4.8 \\
tea-007 & Iron Leaf      & black   & epic      & \$19 & 4.6 \\
tea-008 & Golden Phoenix & special & legendary & \$25 & 5.0 \\
\bottomrule
\end{tabular}
\end{table}

% =========================================================
%  10. DEPLOYMENT & RUNNING
% =========================================================

\section{Deployment and Local Development}

\begin{enumerate}[nosep]
    \item Clone the repository; no dependencies to install.
    \item Open the folder in VS Code and launch \texttt{index.html} with
          the Live Server extension.
    \item For cloud features: create a Supabase project, run
          \texttt{schema.sql} in the SQL Editor, then promote your account
          with
          \texttt{update public.profiles set role = 'admin' where
          email = '<your-email>';}
    \item The production build is a static deploy (currently GitHub
          Pages).
\end{enumerate}

% =========================================================
%  11. LIMITATIONS
% =========================================================

\section{Current Limitations}

\begin{itemize}[nosep]
    \item Payments are simulated; no payment gateway is integrated.
    \item No automated test suite; verification is manual.
    \item Single-file JavaScript engine ($\sim$5.8k lines) is monolithic;
          ES modules would improve maintainability.
    \item Analytics derive from local/cloud order data only --- no event
          tracking.
    \item Images/icons are emoji-based placeholders rather than authored
          pixel art assets.
\end{itemize}

% =========================================================
%  12. FUTURE WORK
% =========================================================

\section{Future Work}

\begin{itemize}[nosep]
    \item Real payment processing (Stripe or similar).
    \item Expanded XP and achievement system with seasonal events.
    \item Interactive tea-world map unlocking region-based blends.
    \item Refactor into ES modules with a lightweight component pattern.
    \item Automated end-to-end tests for checkout and admin flows.
    \item Order confirmation emails via Supabase Edge Functions.
\end{itemize}

% =========================================================
%  13. CONCLUSION
% =========================================================

\section{Conclusion}

TeaQuest demonstrates that personality and engineering rigor can share the
same codebase. In under thirteen thousand lines of dependency-free code it
delivers a complete storefront, a genuine gamification loop, a governed
cloud database and an administrative back office. The architecture's
local-first design keeps the site instant and resilient, while Supabase
row-level security provides defense-in-depth that does not depend on client
code. With payments, modularization and testing on the roadmap, TeaQuest is
well positioned to grow from a polished prototype into a production tea
destination.

\vspace{0.5cm}
\begin{center}
    {\color{teagreen}\rule{0.5\linewidth}{1pt}}\\[6pt]
    \textit{Discover. Brew. Explore.}\\[4pt]
    {\color{teagreen}\rule{6pt}{6pt}}\;
    \textbf{TEAQUEST} --- All pixels reserved.
\end{center}

% =========================================================
%  APPENDIX
% =========================================================

\appendix

\section{localStorage Keys}

\begin{table}[h!]
\centering
\caption{Offline persistence contract.}
\begin{tabular}{@{}lll@{}}
\toprule
\textbf{Key} & \textbf{Type} & \textbf{Contents} \\
\midrule
\texttt{teaquest\_products}    & array  & Product catalog \\
\texttt{teaquest\_users}       & array  & Offline fallback accounts \\
\texttt{teaquest\_cart}        & array  & Cart line items \\
\texttt{teaquest\_favorites}   & array  & Favorited product ids \\
\texttt{teaquest\_orders}      & array  & Placed orders \\
\texttt{teaquest\_currentUser} & object & Active session profile \\
\bottomrule
\end{tabular}
\end{table}

\section{Order Status Workflow}

\begin{lstlisting}[language={},caption={Allowed transitions managed from the admin dashboard.}]
PROCESSING --> SHIPPED --> DELIVERED
     |
     +------> CANCELLED
\end{lstlisting}

\end{document}
